\section{La idea central de Interleaved Thinking}

\subsection{Origen del concepto}

Interleaved Reasoning fue propuesto por primera vez por Apple en un trabajo de 2025: «razonamiento intercalado for modelos de lenguaje grandes mediante aprendizaje por refuerzo». La idea central consiste en descomponer el pensamiento y la respuesta en varios pasos: en cada paso primero se piensa (think) y después se responde (answer), de manera alternada.

\subsection{Ventajas de la descomposición en varios pasos}

Tomemos como ejemplo un problema de razonamiento multisalto: «¿Quién fue el director de la película que ganó el Óscar a la mejor película en el quinto año después de la caída del Muro de Berlín?»

\textbf{Long CoT tradicional}: un razonamiento largo hecho de una sola vez, propenso a errores en los pasos intermedios, que termina produciendo una respuesta incorrecta.

\textbf{Interleaved Thinking}:
\begin{enumerate}
  \item Think 1: se recupera el año de la caída del Muro de Berlín $\to$ 1989
  \item Answer 1: el quinto año = 1994
  \item Think 2: se busca la película que ganó el Óscar a la mejor película en 1994 $\to$ «Forrest Gump»
  \item Answer 2: el director es Robert Zemeckis
\end{enumerate}

\begin{importantbox}{Las tres grandes ventajas de Interleaved Thinking}
\begin{itemize}
  \item \textbf{Latencia del primer token baja}: cada paso think-answer es muy corto, puede tener solo unos cuantos tokens
  \item \textbf{Reward de grano fino}: se puede definir un reward de corrección para la salida de cada paso, lo cual facilita la supervisión de proceso
  \item \textbf{Agentic amigable}: no requiere la participación del usuario; el modelo ejecuta por sí mismo múltiples rondas de tool call
\end{itemize}
\end{importantbox}

\subsection{Orientado a escenarios Agentic}

K2 Thinking puede ejecutar de manera serial entre 200 y 300 pasos de llamadas a herramientas, sin necesidad de participación humana. Desde la perspectiva del message list, se trata de una interacción de una sola ronda entre user y assistant, pero internamente el assistant completa una gran cantidad de ciclos think-act-observe.

\begin{knowledgebox}{Comparación del message list entre Interleaved y Long CoT}
\begin{itemize}
  \item \textbf{Long CoT}: $[\text{user}, \underbrace{\text{<think>}\ldots\text{</think>}}_{\text{cadena de razonamiento muy larga}}, \text{answer}]$
  \item \textbf{Interleaved}: $[\text{user}, \text{think}_1, \text{act}_1, \text{obs}_1, \text{think}_2, \text{act}_2, \text{obs}_2, \ldots, \text{answer}]$
\end{itemize}
\end{knowledgebox}

\subsection{Método de entrenamiento con RL}

El número de pasos de Interleaved Thinking no está fijado de antemano, sino que \textbf{emerge de manera automática} mediante el entrenamiento con RL. El modelo aprende a decidir de manera dinámica, según la complejidad del problema, cuántos ciclos think-answer se necesitan.

\subsection{Resumen de la sección}

Interleaved Thinking es una mejora importante sobre el paradigma Long CoT: descompone la cadena de razonamiento en múltiples pasos think-answer. Reduce la latencia del primer token, admite reward de grano fino y se adapta de forma natural a escenarios de llamadas a herramientas Agentic.
